% =========================================================================
% p2_sec02_multivector_opf.tex — Part II Section II: multi-temporal AC-OPF formulation.
% Security framework cited from Part I rather than repeated.
% =========================================================================
\section{Multi-Temporal AC-OPF with Multi-Vector Coupling}
\label{sec:opf}

The multi-temporal operational campaign exercises the 215-contingency $N-1$ security, two-layer reserve sizing, and bounded-slack framework established in Part~I~\cite{seni142part1}, coupling it directly to multi-vector storage and flexible conversion assets.


\subsection{Multi-temporal AC-OPF with power-to-gas coupling}
\label{sub:opf}

On the converged network model, the benchmark formulates a multi-temporal AC optimal power flow \cite{molzahn2019}. The dispatch prices energy and both reserve directions,
\begin{multline}
\min \sum_{t \in \mathcal{T}} \left[ \sum_{g \in \mathcal{G}} \left( c_g^{P}(P_{g,t}) + c_g^{+} r_{g,t}^{+} + c_g^{-} r_{g,t}^{-} \right) \right. \\
\left. + \sum_{i \in \mathcal{V}} c^{\mathrm{shed}}\, P_{i,t}^{\mathrm{shed}} \right],
\label{eq:opfobj}
\end{multline}
subject to the AC balances, the thermal branch limits, the apparent-power disk of Part~I~\cite{seni142part1}, the hourly ramps $|P_{g,t} - P_{g,t-1}| \le \Delta_g^{\mathrm{ramp}}$, and the storage balances of pumped hydro and hydrogen.

The shedding term is an exact penalty, not a modeling soft spot: the penalty $c^{\mathrm{shed}}$ is chosen strictly above any achievable system marginal cost, $c^{\mathrm{shed}} > \max_{g \in \mathcal{G}} \partial c_g^{P} / \partial P_{g,t}$. At an optimum, KKT stationarity in $P_{i,t}^{\mathrm{shed}}$ reads $c^{\mathrm{shed}} = \lambda_{i,t} + \mu_{i,t}$, with $\lambda_{i,t}$ the multiplier of the nodal balance \eqref{eq:opf_balance} and $\mu_{i,t} \ge 0$ the multiplier of the nonnegativity bound $P_{i,t}^{\mathrm{shed}} \ge 0$. Suppose unserved load were positive at optimality: complementarity would force $\mu_{i,t} = 0$ and hence $\lambda_{i,t} = c^{\mathrm{shed}}$. But whenever the post-event network is statically feasible, a shed-free dispatch exists, and the energy multiplier is then bounded by the marginal cost of the marginal committed unit, $\lambda_{i,t} \le \max_{g} \partial c_g^{P} / \partial P_{g,t} < c^{\mathrm{shed}}$ --- a contradiction. Optimality therefore carries $P_{\mathrm{shed}} = 0$ at every bus and hour in which the post-event network is statically feasible; clause 4 of the PASS battery of Part~I~\cite{seni142part1} is the verified instance of this structural property, not an assumption. The guarantee is quantitative under congestion: with the penalty fixed at the value of lost load, $c^{\mathrm{shed}} = 10{,}000\ \$/\mathrm{MWh}$, the highest registered marginal generation cost of the system below $230\ \$/\mathrm{MWh}$, and worst-case transmission-congestion rents not exceeding $120\ \$/\mathrm{MWh}$, the nodal energy multiplier is bounded by $\lambda_{i,t}^{\max} < 350\ \$/\mathrm{MWh}$; dual feasibility then gives $\mu_{i,t} = c^{\mathrm{shed}} - \lambda_{i,t} > 9{,}650 > 0$ strictly and uniformly under any admissible level of physical congestion, and complementarity forces $P_{i,t}^{\mathrm{shed}} = 0$ identically at every node where the network admits static feasibility.

For each hour $t$ and bus $i$, the equality constraints are the nodal balances, augmented with the power-to-gas coupling terms,
\begin{multline}
\sum_{g \in \mathcal{G}_i} P_{g,i,t} + \sum_{r \in \mathcal{R}_i} P_{r,i,t} + P_{i,t}^{\mathrm{FC}} - P_{i,t}^{\mathrm{EL}} \\
+ P_{i,t}^{\mathrm{shed}} - P_{i,t}^{\mathrm{load}} = P_i(\mathbf{V}_t, \boldsymbol{\theta}_t),
\label{eq:opf_balance}
\end{multline}
with the reactive analogue admitting inverter-coupled absorption/injection within apparent-power capability limits. The intertemporal coupling is imposed by the hydrogen buffer dynamics,
\begin{multline}
E_{h,t+1} = E_{h,t} + \left(\eta_{\mathrm{EL}}\, P_{h,t}^{\mathrm{EL}} - \frac{P_{h,t}^{\mathrm{FC}}}{\eta_{\mathrm{FC}}}\right) \Delta t, \\
E_h^{\min} \le E_{h,t} \le E_h^{\max},
\label{eq:h2_storage}
\end{multline}
where $\eta_{\mathrm{EL}}$ and $\eta_{\mathrm{FC}}$ are the electrolysis and reconversion efficiencies: surplus non-dispatchable renewable generation becomes intertemporal energy inventory, absorbed as flexible electrolytic load in surplus hours and re-injected by the fuel cell in deficit hours. Pumped storage obeys the reservoir analogue with independent pump/turbine efficiencies, bounding the reversible 764~MW fleet of Section~\ref{sub:phes}. The complete problem is nonconvex because the AC equalities are nonconvex. It is solved by a cascade of thirty-two Newton-type configurations spanning Newton--Raphson \cite{tinney1967}, Iwamoto optimal-multiplier step control, fast-decoupled BX iterations \cite{stott1974}, and a Gauss--Seidel fallback, so that no single algebraic pathology can stall the sweep. The Karush--Kuhn--Tucker (KKT) system certifies local stationarity only. On this nonconvex problem the KKT conditions are necessary conditions at a local minimizer. They are not a certificate of global optimality. Stationarity of the Lagrangian, primal feasibility, dual feasibility, and complementarity are archived with the sealed validation output \cite{nocedal2006}, \cite{dommel1968}, \cite{wood2013}. Figure~\ref{fig:p2g_h2} diagrams the electrochemical coupling.

% ---- Figure 2 injection block (verbatim from the sealed master) ----
% Graphic source: sec02_modelado_matematico_ac_h2_sgen.tex, figure environment
% at L63--88 (styles: bloque/flujo/flujoh2/nota; palette seni-*).
% The tikzpicture is extracted verbatim from that source at injection time.
\begin{figure*}[!t]
\centering
\resizebox{0.85\textwidth}{!}{% Figure 2 tikzpicture — English edition (editorial refinement per
% GUIA_BUENAS_PRACTICAS_LATEX_IEEE: white-filled annotations, no line-through
% text; palette seni-* defined in the host preamble).
\begin{tikzpicture}[
  bloque/.style={rectangle, rounded corners=5pt, draw, thick, align=center, font=\small, text width=3.6cm, minimum height=13mm, inner sep=4pt},
  flujo/.style={-{Stealth[length=3mm]}, very thick},
  flujoh2/.style={-{Stealth[length=2.6mm]}, thick, seni-gold!80!black, dash pattern=on 4pt off 2pt},
  nota/.style={font=\scriptsize, align=center, text=seni-dark!85, fill=white, inner sep=2pt},
]
\node[bloque, fill=seni-blue!88, text=white] (red) at (0, 0) {\textbf{AC GRID}\\[2pt] {\scriptsize nodal AC balance\\$Q$ injection/absorption within\\apparent-power capability limits}};
\node[bloque, fill=seni-gold, text=seni-dark] (el) at (5.8, 1.85) {\textbf{ELECTROLYZER}\\[2pt] {\scriptsize $P^{EL}_{h,t} \ge 0$ $\cdot$ efficiency $\eta_{EL}$}};
\node[bloque, fill=seni-gold!45, draw=seni-gold!80!black, text=seni-dark] (buf) at (11.4, 0) {\textbf{H$_2$ BUFFER}\\[2pt] {\scriptsize inventory $E_{h,t}$\\$E^{\min}_h \le E_{h,t} \le E^{\max}_h$}};
\node[bloque, fill=seni-green!82, text=white] (fc) at (5.8, -1.85) {\textbf{FUEL CELL}\\[2pt] {\scriptsize $P^{FC}_{h,t} \ge 0$ $\cdot$ efficiency $\eta_{FC}$}};
%% --- Electric flows: grid <-> converters ---
\draw[flujo, seni-blue!80] (red.north east) .. controls (1.5,1.85) .. (el.west);
\node[nota] at (2.45, 2.72) {electrolysis: $P^{EL}_{h,t} \ge 0$ (absorbs renewable surplus)};
\draw[flujo, seni-green!70!black] (fc.west) .. controls (1.5,-1.85) .. (red.south east);
\node[nota] at (2.45, -2.72) {reconversion: $P^{FC}_{h,t} \ge 0$ (injects during deficit)};
%% --- Hydrogen flows: converters <-> buffer ---
\draw[flujoh2] (el.east) -- node[nota, above, pos=0.55] {$\eta_{EL}\, P^{EL}_{h,t}\, \Delta t$} (buf.north west);
\draw[flujoh2] (buf.south west) -- node[nota, below, pos=0.55] {$P^{FC}_{h,t}\, \Delta t\, /\, \eta_{FC}$} (fc.east);
%% --- Intertemporal coupling ---
\node[nota, text=seni-gold!60!black] at (11.4, -1.35) {intertemporal coupling $\Delta t$};
\end{tikzpicture}
} % extracted verbatim, sec02 L65--85
\caption{Electrochemical power-to-gas coupling of the green-hydrogen vector with the AC network. The electrolyzer absorbs renewable surplus as flexible load ($P^{\mathrm{EL}}_{h,t} \ge 0$, efficiency $\eta_{\mathrm{EL}}$); the storage buffer imposes the intertemporal inventory $E_{h,t+1} = E_{h,t} + (\eta_{\mathrm{EL}} P^{\mathrm{EL}}_{h,t} - P^{\mathrm{FC}}_{h,t}/\eta_{\mathrm{FC}})\,\Delta t$ with box constraints; the fuel cell re-injects power during deficit ($P^{\mathrm{FC}}_{h,t} \ge 0$, efficiency $\eta_{\mathrm{FC}}$). Both convertor terms enter the nodal AC balance. Electric flows in solid strokes; hydrogen flows in dashed gold.}
\label{fig:p2g_h2}
\end{figure*}
